% Job posting: https://ca.linkedin.com/jobs/view/engineer-ai-software-machine-learning-c%2B%2B-at-qualcomm-4457950472
\documentclass[11pt]{article}
\usepackage[left=1in,top=1in,right=1in,bottom=1in]{geometry}
\usepackage[hidelinks]{hyperref}
\usepackage{setspace}
\usepackage[T1]{fontenc}
\usepackage{newtxtext,newtxmath}
\ifdefined\pdfgentounicode
  \input{glyphtounicode}
  \pdfgentounicode=1
\fi
\setstretch{1.1}
\pagenumbering{gobble}
\begin{document}
\pagestyle{empty}

\noindent Nishal Stanislaus Silva, PhD \\
Toronto, ON \\
nishal.silva@hotmail.com \,|\, +1 613 200 2030 \\
nishal.xyz \,|\, linkedin.com/in/nishal-silva \,|\, github.com/ns2max

\vspace{1.5em}

\noindent Dear Hiring Manager,

\vspace{1em}

\noindent I am writing to apply for the Engineer, AI Software/Machine Learning (C++) role with Qualcomm's Machine Learning Engineering group in Markham. The posting is explicit that this is a C++ engineering role, not model building or training, and that is exactly the work I have spent my career doing: writing modern C++17 software that makes trained models run fast and reliably on constrained hardware. My own open-source project, Nebula, is a C++17 audio-feature library with per-feature latency benchmarks on ARM, and my published research detector runs at 14 ms on a Raspberry Pi 4, 74x faster than a DTW baseline, at F1 0.76 (JAES 2026).

\vspace{1em}

\noindent Against the specifics of the posting: I have built and shipped C++/Python real-time inference engines on embedded and IoT hardware across manufacturing sites for MAS Holdings, including reverse-engineering PLC trigger points and building three closed-loop automation subsystems in custom C++, and I design frozen-protocol benchmark and ablation suites as part of my research practice, the same discipline a model-analysis tool for customers requires. My embedded-toolchain experience is concentrated on Raspberry Pi and embedded Linux (Elk Audio OS) rather than Android or QNX specifically, and I have not worked directly with the Hexagon DSP SDK; both are a natural extension of the cross-platform, memory- and latency-constrained C++ work I already do, and I would expect to be productive on them quickly. I use TensorFlow, PyTorch and ONNX routinely as the consumer of the models my tooling has to serve, and I work daily in Git and Docker.

\vspace{1em}

\noindent Qualcomm's Hexagon-based on-device AI stack is the kind of hard compute-and-power-budget problem I have chosen to specialize in rather than one I am pivoting toward, and the Markham site puts that work close to home. I would be glad to bring my benchmark-driven approach to latency and CPU trade-offs to the team's model-analysis tooling from day one.

\vspace{1em}

\noindent I am a Canadian Permanent Resident and do not require sponsorship. I am based in Toronto, a short commute from Markham, and available immediately. I would welcome the chance to discuss how my embedded C++ and ML-systems background can contribute to your team.

\vspace{1.5em}

\noindent Sincerely, \\
Nishal Stanislaus Silva

\end{document}
